\section{A.1. Polynomial categories and strict polynomial functors}
\label{sec:A1}

We distinguish three mathematical objects throughout this section. A
polynomial law is a morphism between modules with scalar-extension data.
The category of divided powers has divided powers of \emph{linear Hom
modules} as its morphisms. A strict polynomial functor is a linear
representation of that latter category. These distinctions refute L1 as
written. They do not diminish the relevance of the manuscript's question
about polynomial laws on Hom modules: that question supplies the precise
bridge between the constructions.

\subsection{The standard module construction, over an arbitrary ring}

\paragraph{A1.1 (Known: polynomial-law representability).}
Let $k$ be any commutative ring, including rings with torsion, and let $M,N$
be arbitrary $k$-modules. A homogeneous polynomial law $f:M\rightsquigarrow
N$ of degree $d\geq0$ means a family of maps
\[
 f_R:R\otimes_k M\longrightarrow R\otimes_k N
 \qquad(R\in\operatorname{CAlg}_k)
\]
natural in the commutative $k$-algebra $R$, with $f_R(rx)=r^d f_R(x)$.
It is not merely its function $f_k$ on $k$-points. There is a natural
isomorphism of $k$-modules
\[
 \operatorname{Pol}^d_k(M,N)
 \simeq\operatorname{Hom}_k(\Gamma^d_k M,N).
 \tag{A1.1}
\]
Here $\Gamma^d$ is the degree-$d$ part of the universal divided-power
algebra; $\Gamma^0 M=k$, $\Gamma^1M=M$. The isomorphism sends $u$ to
$(1\otimes u)\gamma_{d,R}$. Divided powers commute with arbitrary scalar
extension and satisfy
\[
 \Gamma^d(M\oplus N)\simeq
 \bigoplus_{a+b=d}\Gamma^a M\otimes_k\Gamma^b N.
 \tag{A1.2}
\]
The exact references are Roby, Theorem IV.1, pp.~266--267,
Theorem III.3, p.~262, and Theorem III.4, pp.~262--263
\cite{Roby1963}. We opened these statements and their surrounding
definitions. No projectivity or characteristic restriction occurs in
these results. For nonflat modules one must use the universal
divided-power construction, not silently replace it by invariant tensors.

\paragraph{A1.2 (Proved: the full category of bounded polynomial laws).}
For arbitrary $k$, define $\mathcal P_k$ to have $k$-modules as objects and
\[
 \mathcal P_k(X,Y)=\bigoplus_{d\geq0}
     \operatorname{Hom}_k(\Gamma^d X,Y)
\]
as morphisms. These are the polynomial laws with a globally finite degree
bound. Composition is composition of laws after every scalar extension.
This is a strict polynomial category in the representability sense of
cote~137, scans~75--84. Its additive degree-one subcategory is
$\operatorname{Mod}_k$, its finite products are direct sums, and its
representing tensor product and divided powers are the usual ones.

For positive integers or zero $p,q$, let
\[
 \alpha_{q,p}^X:\Gamma^{pq}X\longrightarrow
                    \Gamma^q(\Gamma^pX)
\]
be the linear map representing the degree-$pq$ law
$x\mapsto\gamma_q(\gamma_p(x))$. For $p=0$ this uses the law with
constant value $1\in\Gamma^0X=k$, so $\alpha_{q,0}(1)=\gamma_q(1)$.
For homogeneous $f:\Gamma^pX\to Y$ and
$g:\Gamma^qY\to Z$, composition is exactly
\[
 g\mathbin{\star}f=g\,\Gamma^q(f)\,\alpha_{q,p}^X:
             \Gamma^{pq}X\longrightarrow Z.
 \tag{A1.3}
\]
For an inhomogeneous $f=\sum_{i=0}^m f_i$ and homogeneous $g_q$, its
component of degree $n$ is
\[
 \sum_{\substack{\nu_0+\cdots+\nu_m=q\\
                  \sum i\nu_i=n}}
 g_q\,\mu_\nu^Y
 \left(\bigotimes_{i=0}^m\Gamma^{\nu_i}(f_i)\right)
 \left(\bigotimes_{i=0}^m\alpha_{\nu_i,i}^X\right)
 \Delta_{(i\nu_i)_i}^X .
 \tag{A1.4}
\]
In this formula $\Delta$ is the component of the divided-power
comultiplication
$\Gamma^nX\to\bigotimes_i\Gamma^{i\nu_i}X$, and
$\mu_\nu^Y:\bigotimes_i\Gamma^{\nu_i}Y\to\Gamma^qY$ is multiplication
in $\Gamma^*Y$. In particular every displayed arrow is linear in
$\operatorname{Mod}_k$. The sum is finite, even when constants occur.

\begin{proof}
Under A1.1 a finite sum determines a law, and its degree decomposition is
unique: evaluate it at $tx$ in $N\otimes_k R[t]$ and compare powers of
the indeterminate $t$. Composition of two bounded laws has a bounded
degree, with bound the product of positive degree bounds; the bound zero
case is constant. Its associativity and identities are equality of
functions after every scalar extension, so hold without a faithfulness
assumption about evaluation on $k$-points.

Formula (A1.3) follows by evaluating the represented laws on the universal
element. For (A1.4), expand
$\gamma_q(\sum_i f_i(x))$ using the divided-power sum formula. Each
factor $\gamma_{\nu_i}(f_i(x))$ is represented by
$\Gamma^{\nu_i}(f_i)\alpha_{\nu_i,i}$; combining the factors uses
$\Delta$ in the source and $\mu$ in the target. A1.1 makes equality of
the resulting laws equality of the displayed representing maps.

The law $X\rightsquigarrow Y\oplus Z$ is a pair of laws, of bounded
degree if and only if both components are of bounded degree. Thus
$Y\oplus Z$ is the product and zero is terminal. Addition of laws is
pointwise addition in the target. It is additive in the outer morphism
under composition, while the inner argument is generally nonlinear.
Degree-one laws are precisely linear maps, including in characteristic
two. By (A1.2), a law on $X_1\oplus\cdots\oplus X_r$ of multidegree
$\nu$ is represented by $\bigotimes_i\Gamma^{\nu_i}X_i$.
Consequently the strict tensor comparison in scan~79 is the actual
identity of representing objects
\[
 (X\otimes Y)\otimes\Gamma^a Z
       \simeq X\otimes Y\otimes\Gamma^a Z,
\]
and the analogous statement with any number of factors holds by tensor
associativity. The degree-zero component is $\operatorname{Hom}_k(k,Y)$,
independent of the source. The degree, permutation, scalar, and unit
axioms all follow from these representations and the laws on every
$R$-point. This checks the whole construction, including mixed degrees.
\end{proof}

\paragraph{A1.3 (Refuted: L1; correct relation proved).}
Let $\mathsf P_k$ be the category of finitely generated projective
$k$-modules. The divided-power category used for strict polynomial
functors has the same objects and
\[
 \operatorname{Hom}_{\Gamma^d\mathsf P_k}(V,W)
       =\Gamma^d\operatorname{Hom}_k(V,W).
 \tag{A1.5}
\]
Its composition preserves the fixed degree $d$ and is bilinear. Its
linear representations $\Gamma^d\mathsf P_k\to\operatorname{Mod}_k$
are strict polynomial functors of degree $d$. This is Krause,
\S2, especially the definition on p.~3 and Remark~2.2 on pp.~4--5
\cite{Krause2013}; over a field it is equivalent to Friedlander--Suslin,
Definition~2.1 and Lemma~2.2, \S2 \cite{FriedlanderSuslin1997}.

It is false that the manuscript's $\mathcal P_d(V,W)$ is (A1.5).
The correct relation is that a degree-$d$ polynomial-law action on
\emph{linear Hom modules}, compatible with identities and composition
as laws, is represented by (A1.5), and therefore defines such a linear
representation. The degree-$p$ divided-power functor is one example,
as proved in A1.4. The category $\mathcal P_k$ itself is the category of
bounded polynomial laws, not a fixed-degree divided-power category and
not the category of its representations.

\begin{proof}
Take $k=\mathbb F_2$, $V=k$, $W=k^2$, and $d=2$. Then
$\operatorname{Hom}_k(\Gamma^2 V,W)$ has dimension $2$, while
$\Gamma^2\operatorname{Hom}_k(V,W)$ has dimension $3$: their cardinalities
are $4$ and $8$. The finite certificate records these dimensions and
cardinalities. There is also a structural obstruction: two degree-$d$
morphisms in the manuscript compose in degree $d^2$, so for $d>1$ the
single homogeneous piece is not even a category with the degree-one
identities. There is also an obstruction to an ordinary categorical
equivalence, not merely to an identity on the displayed Hom modules.
For nonzero $k$, the full polynomial-law category has terminal object
$0$ but has no initial object: for every $X$ the distinct constant
laws $0,1:X\rightsquigarrow k$ prevent $\mathcal P_k(X,k)$ from
being a singleton. This remains true on finite projective objects.
For $d>0$, $\Gamma^d\mathsf P_k$ has the zero object $0$, while
for $d=0$ all its Hom sets equal $k$, so it has no terminal object.
Neither can be equivalent to the full polynomial-law category.
The representation category is abelian and has a zero object, giving
the same obstruction to identifying it with $\mathcal P_k$.

For the replacement, apply A1.1 to $M=\operatorname{Hom}_k(V,W)$ and
$N=\operatorname{Hom}_k(FV,FW)$. The degree-$d$ law giving the action of
$F$ is exactly a linear map from (A1.5) to $N$. The identity and
composition diagrams for these linear maps are the representing forms
of the corresponding diagrams of polynomial laws. Thus either set of
data determines the other, with natural transformations characterized
by the same representing diagrams. This proves the asserted bridge;
the equivalence with the standard name is the cited definition/theorem.
\end{proof}

\paragraph{A1.4 (Proved: the closing Hom-law question for modules).}
For all $k$-modules $X,Y$ and every $p\geq0$, the operation
\[
 f\longmapsto\Gamma^p(f):
 \operatorname{Hom}_k(X,Y)\rightsquigarrow
 \operatorname{Hom}_k(\Gamma^pX,\Gamma^pY)
 \tag{A1.6}
\]
is a canonical homogeneous polynomial law of degree $p$. It is natural
in $X,Y$, preserves identities, and is compatible with composition as
an identity of laws. This answers scan~86 affirmatively for its principal
example over \emph{every} commutative ring, without finite generation
or a base-change isomorphism for Hom.

\begin{proof}
For a finite list $f_1,\ldots,f_r\in\operatorname{Hom}_k(X,Y)$, and
$a=(a_1,\ldots,a_r)$ with $|a|=p$, define
\[
 T_a(f_1,\ldots,f_r)=
 \mu_a^Y\left(\bigotimes_j\Gamma^{a_j}(f_j)\right)\Delta_a^X
 \in\operatorname{Hom}_k(\Gamma^pX,\Gamma^pY).
\]
For $h=\sum_j r_j\otimes f_j\in R\otimes_k\operatorname{Hom}_k(X,Y)$
set
\[
 L_R(h)=\sum_{|a|=p}r_1^{a_1}\cdots r_r^{a_r}\otimes T_a(f_1,\ldots,f_r).
 \tag{A1.7}
\]
We verify that this formula is independent of the presentation of $h$.
Permuting terms permutes the indices. Replacing a term by two terms with
the same coefficient is exactly the divided-power sum relation, through
the coassociativity of $\Delta$ and associativity of $\mu$. Combining
two coefficients multiplying the same $f$ is the binomial identity
\[
 \mu_{a,b}\,\Delta_{a,b}=\binom{a+b}{a}\operatorname{id}
       \quad\hbox{on }\Gamma^{a+b},
\]
which follows on divided-power monomials by the ordinary binomial
theorem, or directly by expanding $\gamma_{a+b}((s+t)x)$ over
$k[s,t]$. Moving a scalar of $k$ between $r_j$ and $f_j$ uses
$\Gamma^a(\lambda f)=\lambda^a\Gamma^a(f)$, which follows from
$\gamma_a(\lambda x)=\lambda^a\gamma_a(x)$. These are the defining
relations in a tensor product, together with insertion of zero terms.
Thus (A1.7) is a well-defined natural family, homogeneous of degree $p$.
For $R=k$ and $r=1,r_1=1$ it is $\Gamma^p(f)$.

Compatibility with composition and identities can be checked with
independent indeterminates for the finitely many coefficients of both
maps. Over $k[\mathbf s,\mathbf t]$, (A1.7) is the coefficient expansion
of the ordinary functorial map $\Gamma^p(\sum t_i f_i)$, using A1.1's
scalar-extension isomorphism for divided powers. Equality follows from
functoriality of $\Gamma^p$ there. Coefficient equality occurs in
$\operatorname{Hom}_k(\Gamma^pX,\Gamma^pY)[\mathbf s,\mathbf t]$:
its map into maps to $(\Gamma^pY)[\mathbf s,\mathbf t]$ is injective,
because polynomial coefficients can be read separately. Specialization
to any $R$ proves the law identity. We have nowhere assumed
$R\otimes\operatorname{Hom}(X,Y)\simeq
\operatorname{Hom}_R(R\otimes X,R\otimes Y)$.
\end{proof}

\subsection{What the abstract axioms do and do not imply}

For precision, call the data explicitly listed on scans~82--83 the
\emph{bare tensor/divided-power data}: a $k$-linear additive symmetric
monoidal category $\mathcal C$, a coherent exponential functor
$G^*:(\mathcal C,\oplus)\to(\operatorname{Gr}_{\geq0}\mathcal C,\otimes)$,
$G^0=\mathbf1$, $G^1=\operatorname{Id}$, the stipulated scalar degrees,
and natural maps
\[
 \alpha_{p,q}:G^{pq}\longrightarrow G^pG^q
\]
satisfying the typed associative diagram
\[
 G^p(\alpha_{q,r})\alpha_{p,qr}
       =\alpha_{p,q}^{G^r}\alpha_{pq,r}.
 \tag{A1.8}
\]
The units are $\alpha_{1,q}=\alpha_{q,1}=\operatorname{id}$,
$\alpha_{0,q}=\operatorname{id}_{\mathbf1}$, and, when included,
$\alpha_{q,0}:\mathbf1\overset\sim\longrightarrow G^q\mathbf1$.
The superscript in the top-right arrow of the scan~81 diagram is doubtful;
(A1.8) is its unique type-correct reading. A reading with a different
source is not a mathematical composition. Our counterexamples concern
this strongest type-correct homogeneous reading.

\paragraph{A1.5 (Proved: independence of the omitted unit axiom).}
The condition $\alpha_{q,0}$ being an isomorphism is independent of all
the other bare tensor/divided-power axioms, even with exponential
coherence, scalar degrees, and the complete homogeneous associative
diagram.

\begin{proof}
Take $k=\mathbb F_2$, $\mathcal C=\mathsf{Vect}^{\mathrm{fd}}_k$ with
its ordinary tensor symmetry, and $G^n(V)=\Lambda^nV$. Exterior powers
are exponential; in characteristic two the usual signs in the
exponential symmetry disappear. We have $G^0=k$, $G^1=\operatorname{Id}$
and $G^n(\lambda f)=\lambda^nG^n(f)$ for both scalars $0,1$.

For $p,q>0$, define $\alpha_{p,q}$ on an exterior product of $pq$
vectors by summing over unordered partitions into $p$ blocks of size
$q$, taking the exterior product inside each block and then the exterior
product of the $p$ resulting elements of $\Lambda^qV$. This expression
is multilinear. It vanishes when two input vectors coincide: terms with
both in one block vanish; the transposition of their positions pairs
the remaining terms, which cancel in characteristic two. If $q=1$, the
fixed partition instead gives repeated outer factors and vanishes
directly. Thus it descends to $\Lambda^{pq}V$, and is natural in $V$.

Both sides of (A1.8) sum over precisely the same three-level partitions:
$p$ groups, each containing $q$ blocks, each containing $r$ labels.
There is a bijection between first choosing the top groups and refining
them, and first choosing the bottom blocks and grouping them. Each
nested partition occurs once, proving equality in every dimension.
The maps for $p=1$ or $q=1$ are identities. Set $\alpha_{0,q}$ to the
identity of $k$, and $\alpha_{p,0}:k\to\Lambda^p k$ to the canonical
map, the identity for $p=0,1$ and zero for $p\geq2$. Every zero-index
instance of (A1.8) either is an identity of $k$ or has both sides zero;
this verifies the remaining cases.

Finally $G^2(k)=\Lambda^2k=0$, so $k\to G^2(k)$ is not an
isomorphism. This is the required model, not just the observation that
the axiom was not written. The program checks the finite unit witness
and all nested partition identities with $pqr\leq8$ (38 triples);
the partition proof establishes all degrees.
\end{proof}

\paragraph{A1.6 (Refuted: the printed full equivalence).}
The bare tensor/divided-power data, \emph{including} the invertible
zero-index units, do not imply a strict polynomial category with the
specified exponential and tensor-unit comparison maps. Consequently
``the two notions are equivalent'' on scan~84 is false for exactly the
data listed on scans~82--83. The homogeneous assertion alone is true.

\begin{proof}
Take finite-dimensional $\mathbb F_3$-vector spaces, the standard tensor
product, and the standard divided powers with their standard exponential
comparison maps. Put $t_2=-1$ and $t_n=1$ for every $n\ne2$, including
$0,1$. Replace the standard substitution maps by
\[
 \alpha'_{p,q}=c_{p,q}\alpha_{p,q},\qquad
 c_{p,q}=\frac{t_{pq}}{t_p t_q^p}\in\mathbb F_3^\times.
\]
Naturality and the exponential data are unchanged. The scalar identity
\[
 c_{p,qr}c_{q,r}^p=c_{pq,r}c_{p,q}
\]
follows by cancellation of the displayed $t$'s, so (A1.8) holds. The
one-index units and $\alpha'_{0,q}$ are identities, and every
$\alpha'_{q,0}$ is an isomorphism. Thus every stipulated axiom holds.
But $\alpha'_{2,0}=-\alpha_{2,0}$.

Suppose a strict polynomial category inducing these data existed.
Let $U=k=\mathbf1$, let $e_0:0\to U$ be the universal constant
represented by $1\in\operatorname{Hom}_k(k,k)$, and let
$m:U\times U\to U$ be the universal bilinear map, represented by the
standard projection
$\Gamma^2(U\oplus U)\to U\otimes U\simeq U$.
The strict mixed comparison of multidegree $(0,1)$, with the fixed
tensor-unit comparison, gives $m(e_0,\operatorname{id}_U)
=\operatorname{id}_U$. Composing this equality with $e_0$ yields
\[
                 m(e_0,e_0)=e_0.
 \tag{A1.9}
\]
On the other hand the homogeneous formula (scan~83), applied to the
degree-zero pair $(e_0,e_0)$ and the degree-two map $m$, gives
\[
 m(e_0,e_0)=m\,\Gamma^2((1,1))\,\alpha'_{2,0}=-e_0,
\]
since the mixed coefficient of $\gamma_2(1,1)$ is $1$. This contradicts
(A1.9) over $\mathbb F_3$.

This is a statement about equivalence of the \emph{structured} data in
the manuscript. Rescaling homogeneous representatives cannot remove
the counterexample while keeping the exponential comparisons: a
monoidal rescaling would require $t_{a+b}=t_a t_b$, whereas here
$t_2\ne t_1^2$. The program verifies the unit contradiction and the
scalar associativity for all $0\leq p,q,r\leq16$; the cancellation
proof establishes it for all indices.
\end{proof}

\paragraph{A1.7 (Proved: an exact repair beyond homogeneous degrees).}
There is a necessary and sufficient coherence condition for the desired
extension, with no assumption on $k$ or the existence of infinite sums
in $\mathcal C$. We give it explicitly. Write
\[
 B_mX=\bigoplus_{i=0}^m G^iX.
\]
For $d,m\geq0$ define a natural map
\[
 D_{d,m}:B_{dm}X\longrightarrow B_d(B_mX)
 \tag{A1.10}
\]
by the following components. In $G^q(B_mX)$, for $0\leq q\leq d$,
use the exponential decomposition indexed by
$\nu=(\nu_0,\ldots,\nu_m)$, $\sum\nu_i=q$. The component with
weight $n=\sum i\nu_i$ is zero on other $G^aX$, and on $G^nX$ is
\[
 G^nX\xrightarrow{\Delta_{(i\nu_i)}}
 \bigotimes_iG^{i\nu_i}X
 \xrightarrow{\bigotimes_i\alpha_{\nu_i,i}}
 \bigotimes_iG^{\nu_i}(G^iX).
 \tag{A1.11}
\]
Here $\Delta$ is obtained by applying $G^n$ to the additive diagonal
and projecting to the specified exponential component. All maps in
(A1.10)--(A1.11) belong to $\mathcal C$.

The maps $D_{d,m}$ must satisfy the two counit diagrams, explicitly
\[
 \epsilon_{B_mX}D_{1,m}=\operatorname{id}_{B_mX},\qquad
 B_d(\epsilon_X)D_{d,1}=\operatorname{id}_{B_dX},
\]
where $\epsilon_X:B_1X\to X$ is projection to degree one, and the
following finite coassociativity equations for all $a,b,c\geq0$:
\[
 B_a(D_{b,c})D_{a,bc}
       =D_{a,b}^{B_cX}D_{ab,c}:
 B_{abc}X\longrightarrow B_aB_bB_cX.
 \tag{A1.12}
\]
The projection compatibilities in the bounds $d,m$ are also required.
These equations say that the pro-endofunctor $B=(B_m)_{m\geq0}$ is a
comonad with this specified $D$ and $\epsilon$. We impose the following
two explicit grading and representability coherences as well; no
minimality or redundancy assertion about this list is needed.

First, each $D_{d,m}$ must preserve the multigradings obtained by
replacing $X$ with a finite sum $\bigoplus_jX_j$. To make this a
literal matrix condition, expand its target twice with the exponential
isomorphisms. An inner summand has the form
$U_\beta=\bigotimes_jG^{\beta_j}X_j$, with $|\beta|\leq m$.
A target summand is then $\bigotimes_\beta G^{\nu_\beta}U_\beta$,
with $\sum\nu_\beta\leq d$, and has weight
$w_j=\sum_\beta\nu_\beta\beta_j$. The component from the source
multidegree $\lambda$ to this summand must vanish unless $\lambda=w$.
All sums and all tested components are finite.

Second, for every $p,q\geq0$ and $X_1,\ldots,X_q$, require the map
identity
\[
 \operatorname{pr}_{(1,\ldots,1)}\,
 G^q\bigl((G^p\operatorname{pr}_j)_j\bigr)\,
 \alpha_{q,p}^{\oplus_jX_j}
       =\operatorname{pr}_{(p,\ldots,p)}.
 \tag{A1.12N}
\]
Both sides are maps from $G^{pq}(\bigoplus_jX_j)$ to
$\bigotimes_jG^pX_j$. The tuple in the middle is the linear map
$G^p(\bigoplus_jX_j)\to\bigoplus_jG^pX_j$; the outer projection
uses the exponential comparison. When $p=0$ this tuple is the
diagonal of the unit. Empty tensor products and the $q=0$ case use
the stated unit comparisons. Thus (A1.12N) forces both the identity
on the desired component and zero on every other component. This
last clause is essential over finite fields, where scalar eigenvalues
alone do not exclude off-diagonal components.

This is a family of finite, typed diagrams in the original category,
not a demand for an infinite divided-power object. Bare homogeneous
associativity (A1.8) is only the part of (A1.12) with a single
homogeneous component at each stage.

The category of strict polynomial categories with the chosen strict
representing data is equivalent to the category of bare data satisfying
these additional diagrams. Morphisms on both sides are structure
preserving functors. Thus the correction is an exact criterion, not
merely a sufficient special case. The standard module construction
satisfies it.

\begin{proof}
First suppose a strict polynomial category is given. For each $m$ its
universal morphisms $e_i:X\to G^iX$ form a map
$\eta_m:X\to B_mX$ by the product property. Compose $\eta_m$ with
the universal maps out of $B_mX$ of degrees $0,\ldots,d$.
Strict mixed representability identifies the resulting representing
map with (A1.10): on a multidegree $\nu$ its factors are the composites
$e_{\nu_i}\circ e_i$, represented by $\alpha_{\nu_i,i}$; putting the
variables equal uses the displayed $\Delta$. Thus the formula is
forced. Forgetting higher components gives the projection
compatibilities. The identity morphism gives the counits. Associating
three maps $\eta_c,\eta_b,\eta_a$ in the two ways gives (A1.12).
The degree axiom applied with all variables kept separate proves
the asserted zero-component grading conditions. Strict representability
applied to the $q$ universal degree-$p$ maps in disjoint variables
says that their outer universal multilinear product is the universal
multidegree-$(p,\ldots,p)$ map. Writing its representing map by
(A1.3) is exactly (A1.12N), proving necessity of that condition.

Conversely, assume the finite diagrams. Define
\[
 \operatorname{Hom}_{\mathcal P}(X,Y)
   =\varinjlim_m\operatorname{Hom}_{\mathcal C}(B_mX,Y)
   =\bigoplus_i\operatorname{Hom}_{\mathcal C}(G^iX,Y).
\]
For representatives $f:B_mX\to Y$, $g:B_dY\to Z$, set
\[
       g\star f=g\,B_d(f)\,D_{d,m}.
\]
The projection compatibilities make this independent of the bounds.
The two counits make $\epsilon_X$ the identity. Substitution of the
formula for three morphisms, followed by naturality of $D$, reduces
associativity to (A1.12) exactly. This is also the usual elementary
proof of the co-Kleisli construction, performed with finite bounds.

Finite products and addition in the target follow from additivity of
$\mathcal C$. The pure degree-one components of the counits and the
one-index units identify the additive subcategory with $\mathcal C$.
The weighted-zero clauses and the explicit multigrading condition
give the degree rules; the exponential comparisons give the
multigradings and their permutation compatibility. For the strict
representing maps, group the incoming universal maps by their common
degree $p$. On a group of $q$ equal degrees, (A1.11) uses
$\alpha_{q,p}$ and its relevant full component is the identity
projection by (A1.12N). Groups of distinct degrees are combined by
the disjoint-variable components of $\Delta$, which are the
exponential identity comparisons. Their tensor is therefore precisely
the required universal multidegree projection, including its
vanishing on all other components. This proves the strict
$G$-representability comparisons with any extra variables present.
For a strict tensor comparison the tensor variable has degree one;
its insertion uses only the same identity components and the
associativity of the specified tensor comparisons. The case $p=0$
of (A1.12N) gives the prescribed tensor-unit and constant comparisons.
Thus all strictness requirements, including the ones violated by
A1.6, are verified.

Both constructions leave every representing Hom module and every
universal map unchanged. They are inverse by representability, and
the same argument applies to structure preserving functors and their
natural isomorphisms. For actual divided powers, A1.2 identifies the
formula with associative composition of laws, so all diagrams hold.
\end{proof}

\paragraph{A1.8 (Refuted: automatic strict enrichment of bare abstract data).}
Even with $\mathbf1\simeq G^q\mathbf1$ for every $q$, the bare data
do not imply a natural polynomial-law enrichment of
$f\mapsto G^p(f)$ compatible with the functor structure. The
affirmative answer A1.4 therefore cannot be transferred to every
abstract tensor/divided-power category as printed.

\begin{proof}
Let $\mathcal C$ be the category of finite-dimensional nonnegatively
\emph{internally} graded $\mathbb F_2$-vector spaces and degree-zero
linear maps, with its usual graded tensor product. Signs are trivial
in characteristic two. Let $T$ discard internal degree zero, and let
$D^n$ denote standard divided powers in this category. Define an
\emph{external} grading by
\[
 G^nX=\bigoplus_{a+3b=n}D^aX\otimes D^b(TX).
 \tag{A1.13}
\]
These are finite objects. Since $T$ is additive, (A1.13) is a coherent
exponential functor. It has $G^0=\mathbf1$, $G^1=\operatorname{Id}$,
and $G^n\mathbf1=\mathbf1$ for every $n$. The scalar-degree condition
holds for $0,1\in\mathbb F_2$. There are natural exponential maps
$i_n:D^n\to G^n$ and $\pi_n:G^n\to D^n$ from the summand $b=0$,
with $\pi_n i_n=\operatorname{id}$.

For $p,q\geq2$ put
\[
 \alpha^G_{p,q}=
 i_p^{G^qX}\,D^p(i_q^X)\,\alpha^D_{p,q}\,\pi_{pq}^X.
 \tag{A1.14}
\]
For a one-index use the identity, and for a zero-index use the
canonical unit maps. All types in (A1.14) are the successive objects
$G^{pq}X,D^{pq}X,D^pD^qX,D^pG^qX,G^pG^qX$.
For $p,q,r\geq2$, substitution in (A1.8), naturality of $i,\pi$,
and cancellation of each $\pi_j i_j$ reduce both sides to the same
threefold standard divided-power substitution, preceded by
$\pi_{pqr}$ and followed by the corresponding inclusions. This proves
associativity. One-index cases are identities. If an index is zero,
all relevant values at $\mathbf1$ are $\mathbf1$, and the remaining
diagram is the identity unit diagram. Thus every bare axiom, including
the omitted isomorphism axiom, holds.

Now put a vector space $V$ entirely in internal degree one. Then
\[
 G^3V=D^3V\oplus V,
\]
with the two summands in internal degrees three and one. On the second
summand $G^3(f)$ is $f$ itself. Suppose there were a natural homogeneous
cubic law on Hom modules extending this functor. Projecting to that
internal-degree-one block would give laws
\[
 \ell_{V,W}:\operatorname{Hom}(V,W)\rightsquigarrow
             \operatorname{Hom}(V,W)
\]
of degree three, natural for linear pre- and postcomposition, with
$\ell_k$ equal to the identity on $k$-points. For $V=W=k$,
homogeneity forces $\ell_R(r)=r^3$ for every $k$-algebra $R$.
Coordinate inclusions and projections then force $\ell_R$ on matrices
to be entrywise cubing. Naturality for the linear map
$k^2\to k$, $(x,y)\mapsto x+y$, would require
$(x+y)^3=x^3+y^3$ over every $R$. In
$R=\mathbb F_4=\mathbb F_2[\omega]/(\omega^2+\omega+1)$,
take $x=1,y=\omega$. The left side is $1$ and the right side is $0$.
This finite contradiction is checked exhaustively on all 16 pairs.

The conclusion concerns the coherent enrichment asked for in the
strict-polynomial comparison. A separately chosen cubic law on an
individual finite Hom space can be obtained by entrywise cubing in
chosen bases; it fails precisely the naturality just proved necessary.
We do not confuse that weaker existence assertion with functorial
polynomial enrichment.
\end{proof}

\subsection{The grading as actual data}

\paragraph{A1.9 (Proved: the exact scalar ambiguity).}
Suppose an abelian group $A=\bigoplus_{n\geq0}A_n$ has operators
$u(\lambda)$ acting on $A_n$ as $\lambda^n$, for all $\lambda\in\mathbb Z$.
For $n\ne m$, put
\[
 D_{n,m}=\gcd_{\lambda\in\mathbb Z}|\lambda^n-\lambda^m|,
 \qquad 0^0=1.
\]
The scalar eigenspace of weight $n$ is exactly
\[
 \{x:u(\lambda)x=\lambda^n x\ \forall\lambda\}
  =A_n\oplus\bigoplus_{m\ne n}A_m[D_{n,m}].
 \tag{A1.15}
\]
Here $A[D]$ denotes the subgroup killed by $D$. If one index is zero,
$D_{n,m}=1$. For positive distinct indices, set $s=\min(n,m)$ and
$h=|n-m|$. The only prime divisors of $D_{n,m}$ are the primes
$\ell$ with $(\ell-1)\mid h$. More precisely
\[
 v_\ell(D_{n,m})=
 \begin{cases}
 \min(s,1+v_\ell(h)),&\ell>2, (\ell-1)\mid h,\\
 0,&\ell>2, (\ell-1)\nmid h,\\
 \min(s,1),&\ell=2, h\text{ odd},\\
 \min(s,2+v_2(h)),&\ell=2, h\text{ even}.
 \end{cases}
 \tag{A1.16}
\]
Thus inverting the listed primes is sufficient to separate a prescribed
set of weights, but is not necessary: torsion-free groups already have
no ambiguity. The necessity wording in R6 and the study must be
weakened accordingly. At degree zero there is no ambiguity at all.

For a $k$-module with scalar operators for every $\lambda\in k$, the
exact same answer replaces $(D_{n,m})$ by the ideal
\[
 I_{n,m}(k)=(\lambda^n-\lambda^m:\lambda\in k),
\]
and $A_m[D]$ by the submodule annihilated by that ideal. In particular,
over $\mathbb F_q$, positive weights coincide on scalars precisely
when $(q-1)\mid(n-m)$; an infinite field separates all distinct
weights. For multidegrees $\mathbf n,\mathbf m$, the relevant ideal
is $\sum_j I_{n_j,m_j}(k)$. Thus the integral prime-field congruence
must not be asserted for all scalars of a larger finite field.

If the scalar action and grading extend naturally to every
commutative base algebra, with the weight identity there, all weights
are distinguished integrally: the action of the indeterminate $t$
separates them without any prime inversion.

\begin{proof}
An element of a direct sum has finitely many nonzero components.
Comparing those components shows that it belongs to the weight-$n$
eigenspace exactly when each other component is killed by every
$\lambda^n-\lambda^m$, hence by their gcd. This proves (A1.15).
For the zero index use $\lambda=0$.

For positive indices the differences have absolute value
$\lambda^s(\lambda^h-1)$. Modulo $\ell$, the unit group of the prime
field is cyclic of order $\ell-1$. Therefore every difference vanishes
modulo $\ell$ exactly when $(\ell-1)\mid h$. For the valuation, inputs
divisible by $\ell$ have minimal valuation $s$, attained at $\lambda=\ell$.
For odd $\ell$ satisfying the divisibility condition, write
$h=(\ell-1)\ell^v h'$ with $\ell\nmid h'$. An element whose reduction
generates $(\mathbb Z/\ell^2)^\times$ has
$v_\ell(\lambda^{\ell-1}-1)=1$, and the elementary binomial expansion
of $(1+\ell a)^{\ell^v h'}$ has exact valuation $1+v$ when
$\ell\nmid a$. Every unit has at least that valuation, by Fermat's
congruence followed by the same expansion. This proves the odd-prime
case. The cyclic generator modulo $\ell^2$ can also be obtained by
lifting a generator modulo $\ell$: among its $\ell$ lifts, at least
one has nonzero first-order coefficient, since the derivative of
$x^{\ell-1}$ is nonzero modulo $\ell$.

For odd $\lambda$ and odd $h$, the minimum is one, attained at
$\lambda=3$. For even $h$, successive factorization of
$\lambda^h-1$ gives
\[
 v_2(\lambda^h-1)
   =v_2(\lambda-1)+v_2(\lambda+1)+v_2(h)-1.
\]
The sum of the first two valuations is at least three, with equality
at $\lambda=3$, giving the final case of (A1.16). Taking the minimum
with $s$ proves the formula. The finite certificate verifies 72
instances against gcds of the evaluations $0\leq\lambda\leq200$;
the valuation argument establishes all indices.

Finally, over $k[t]$ a putative component of two different weights
would satisfy $(t^n-t^m)x=0$ in $A\otimes_k k[t]$. Its distinct
polynomial coefficients are $x$ and $-x$, so $x=0$, even with torsion.
In characteristic $p$, the familiar equal functions $\lambda$ and
$\lambda^p$ on $\mathbb F_p$ are therefore distinct laws: at $t$ their
values differ. This verifies both the withdrawn uniqueness argument
and the manuscript's proposed all-base-change remedy.

The general-ring formula follows by the same component comparison.
For a finite field, its cyclic multiplicative group gives the stated
criterion; for an infinite field, a nonzero polynomial cannot vanish
at every scalar, so some difference is a unit. In several variables,
setting all but one scalar equal to one shows that the ideal of
differences of the two monomials contains each $I_{n_j,m_j}(k)$.
Conversely telescoping the product difference puts every generator
in their sum, proving equality. Independent indeterminates give the
all-base-change separation in several variables as well.
\end{proof}

\paragraph{Literature scope and attribution.}
We searched exact titles and the terms ``polynomial categories'',
``exponential functors'', ``polynomial laws on Hom'', and ``universal
property of strict polynomial functors''. In addition to the opened
Roby, Friedlander--Suslin and Krause sources, we opened the arXiv
LaTeX source of Touz\'e's July~2026 preprint \cite{Touze2026},
especially Proposition~1, Example~2, Definition~4 and
Theorem~5. Its universal property concerns exact tensor functors from
the \emph{category of strict polynomial functors}; it does not identify
the two different Hom modules in A1.3, nor supply missing coherence
for the bare data. We claim no priority for the constructions or
counterexamples. No point in this section is taken from Michel Hua's
project.
